% ####################################################################
% Basic configuration and packages
% ####################################################################
% Package for discovering wrong and outdated usage of LaTeX.
% More information to be found in l2tabu English version.
\RequirePackage[l2tabu, orthodox]{nag}
% Class of LaTeX document: {size of paper, size of font}[document class]
\documentclass[a4paper,11pt]{article}



% ------------------------------------------------------
% Packages not tied to particular normal language
% ------------------------------------------------------
% This package should improved spaces in the text
\usepackage{microtype}
% Add few important symbols, like text Celcius degree
\usepackage{textcomp}



% ------------------------------------------------------
% Polonization of LaTeX document
% ------------------------------------------------------
% Basic polonization of the text
\usepackage[MeX]{polski}
% Switching on UTF-8 encoding
\usepackage[utf8]{inputenc}
% Adding font Latin Modern
\usepackage{lmodern}
% Package is need for fonts Latin Modern
\usepackage[T1]{fontenc}



% ------------------------------------------------------
% Setting margins
% ------------------------------------------------------
\usepackage[a4paper, total={14cm, 25cm}]{geometry}



% ------------------------------------------------------
% Setting vertical spaces in the text
% ------------------------------------------------------
% Setting space between lines
\renewcommand{\baselinestretch}{1.1}

% Setting space between lines in tables
\renewcommand{\arraystretch}{1.4}



% ------------------------------------------------------
% Packages for scientific papers
% ------------------------------------------------------
% Switching off \lll symbol, that I guess is representing letter "Ł"
% It collide with `amsmath' package's command with the same name
\let\lll\undefined
% Basic package from American Mathematical Society (AMS)
\usepackage[intlimits]{amsmath}
% Equations are numbered separately in every section
\numberwithin{equation}{section}

% Other very useful packages from AMS
\usepackage{amsfonts}
\usepackage{amssymb}
\usepackage{amscd}
\usepackage{amsthm}

% Package with better looking calligraphy fonts
\usepackage{calrsfs}

% Package with better looking greek letters
% Example of use: pi -> \uppi
\usepackage{upgreek}
% Improving look of lambda letter
\let\oldlambda\Lambda
\renewcommand{\lambda}{\uplambda}




% ------------------------------------------------------
% BibLaTeX
% ------------------------------------------------------
% Package biblatex, with biber as its backend, allow us to handle
% bibliography entries that use Unicode symbols outside ASCII
\usepackage[
language=polish,
backend=biber,
style=alphabetic,
url=false,
eprint=true,
]{biblatex}

\addbibresource{Systemy-operacyjne-Bibliography.bib}





% ------------------------------------------------------
% Defining new environments (?)
% ------------------------------------------------------
% Defining enviroment "Wniosek"
\newtheorem{corollary}{Wniosek}
\newtheorem{definition}{Definicja}
\newtheorem{theorem}{Twierdzenie}





% ------------------------------------------------------
% Local packages
% You need to put them in the directory Local-packages
% ------------------------------------------------------
% Local configuration of this particular file
\usepackage{./Local-packages/local-settings}

% Package containing various command useful for working with a text
\usepackage{./Local-packages/general-commands}
% Package containing commands and other code useful for working with
% mathematical text
% \usepackage{math-commands}





% ------------------------------------------------------
% Package "hyperref"
% They advised to put it on the end of preambule
% ------------------------------------------------------
% It allows you to use hyperlinks in the text
\usepackage{hyperref}










% ####################################################################
% Title and author of the text
\title{Historia informatyki \\
  {\Large Błędy i~uwagi}}

\author{Kamil Ziemian \\
  \email}


% \date{}
% ####################################################################










% ####################################################################
% Beginning of the document
\begin{document}
% ####################################################################





% ######################################
% Title of the text
\maketitle
% ######################################





% ######################################
\section{Brian W. Kernighan \textit{Jak Unix tworzył historię},
  \parencite{Kernighan-Jak-Unix-tworzyl-historie-Pub-2021}}

\label{sec:Kernighan-Jak-Unix-tworzył-historię}
% ######################################


% ##################
\CenterBoldFont{Uwagi do~konkretnych stron}

\noindent
\Str{23} Tutaj pierwszy raz napotykamy~się na problem przeliczenia starych
jednostek komputerowych na te używane współcześnie. Ze względu na upływ
czasu nie jest to wcale zadanie proste i~chyba sam Kernighan miał z~tym
problemy. Wedle tego co napisano na tej stronie, symbol $\text{K}$ oznacza
$1 \, 024 = 2^{ 10 }$ jednostek, gdyż wtedy $32\text{K}$ to rzeczywiście
$32 \ 768$ jednostek. Jednak podane tu przeliczenie starych jednostek
na~kilobajty napotyka pewien problem.

Zacznijmy od tego, że~w~naszej ocenie autor przez kilobajt
rozumie~$1 \, 024$ bajtów, zaś~pojedynczy bajt to oczywiście $8$~bitów.
Ile pamięci daje nam $32\text{K}$ słów o~długości $36$~bitów? Ponieważ~$36$
bitów to $4.5$ bajta, więc $32\text{K}$ taki słów, to $4.5 \cdot 32 = 144$
kilobajtów, czyli wyraźnie więcej niż podane tutaj $128$ kilobajtów. Jeśli na
chwilę przyjmiemy, że~autor przez kilobajt rozumie $1 \, 000$ bajtów, to
wówczas mamy
\begin{equation}
  \label{eq:Kernighan-Jak-Unix-etc-01}
  144 \cdot 1 \, 024 = 144 \cdot 1.024 \cdot 1 \, 000 = 147.456 \cdot 1 \, 000,
\end{equation}
znów więcej, niż~$128$ kilobajtów.

Wychodzi więc na to, że~autor pomylił~się, przy przeliczaniu jednych
jednostek na drugie, co nie jest niczym niezwykłym. Choć w~tej książce
wielokrotnie pojawiają~się fragmenty, w~których prezentowane są dane
w~starych jednostkach i~odpowiadające im wartości w~nowych, nie będziemy
w~każdym z~tych przypadków sprawdzać, czy podane wartości zgadzają~się
ze sobą. Poprzestaniemy na stwierdzeniu, że~w~każdej z~tych sytuacji należy
zachować czujność i~ewentualnie przeliczyć te wielkości samemu.

\VerSpaceFour





% \noindent
% \StrWierszeGora{31}{16--17}

% \VerSpaceFour





% \noindent
% \Str{76--77}

% \VerSpaceFour





% \noindent
% \StrWierszeDol{82}{9}

% \VerSpaceFour





% \noindent
% \Str{91} Przedstawiony tutaj kod w~języku~C nie jest zbyt elegancki, ze
% względu na brak odpowiednich wcięć. Takie wygląd jest jednak
% usprawiedliwiony ograniczoną ilością miejsca na stronie, dlatego
% w~przyszłości nie będziemy więcej komentować tego typu kwestii.

% \VerSpaceFour





% \noindent
% \StrWierszDol{127}{17} Angielska nazwa dla procesu z~drugiego planu,n czyli
% \textit{deamon}, zostało tu przetłumaczona jako „demon”. Jednak w~języku
% polskim bardziej poprawny tłumaczeniem tego terminu jest „daimion”.





% ######################################
\newpage

\CenterBoldFont{Błędy}


\begin{center}

  \begin{tabular}{|c|c|c|c|c|}
    \hline
    Strona & \multicolumn{2}{c|}{Wiersz} & Jest
    & Powinno być \\ \cline{2-3}
           & Od góry & Od dołu & & \\
    \hline
    \hphantom{0}61 & \hphantom{0}9 & & instrukcji & instrukcji obsługi \\
    \hphantom{0}78 & & \hphantom{0}5 & funkcjonalne & funkcyjne \\
    101 & 15 & & 1976 & 1966 \\
    % & & & & \\
    % & & & & \\
    % & & & & \\
    \hline
  \end{tabular}

\end{center}

\VerSpaceSix


\noindent
\StrWierszeGora{54}{7--8} \\
\Jest doktoratu \\
<< \\
\PowinnoByc doktoratu<< \\
\StrWierszeDol{79}{5} \\
\Jest obsługi skomplikowanego \\
\PowinnoByc obsługi \\



% ######################################










% ####################################################################
% ####################################################################
% Bibliography

\printbibliography





% ############################
% End of the document

\end{document}
